\section{The \texorpdfstring{$p$}{p}-part in rank at most one}
\label{sec:t-cascade}

In analytic rank $0$ and $1$, the recognition predicate contains the
scalar Strong BSD identity itself as a hypothesis
(\Cref{def:closure:pi-bsd}).  This section records, for context, the
form in which known results are expected to bear on that hypothesis: they
give the $p$-part of the identity for suitable primes $p$.  As with the
other inputs, we state precisely what is assumed; the main theorem does not
use this section.

\begin{definition}[Rank-at-most-one input]\label{def:closure:t-cascade-import}
The rank-at-most-one input is a list of forty propositions, each with an
asserted proof, standing for published results in two branches: for curves
with complex multiplication, Rubin's Iwasawa theory over imaginary
quadratic fields~\citep{Rubin1991}, the main conjecture at supersingular
primes of Pollack--Rubin~\citep[Theorem~7.3]{PollackRubin2004}, and work of
Burungale--Castella--Skinner--Tian~\citep{BurungaleCastellaSkinnerTian2022}
at good ordinary primes, together with the archimedean, Heegner-point and
uniformization ingredients needed to compare $p$-adic and classical
quantities; for curves without complex multiplication, the work of
Skinner--Urban~\citep{SkinnerUrban2014}, Wei Zhang~\citep{WeiZhang2014},
Castella et al.~\citep{CHKLL2025} and Yan--Zhu~\citep{YanZhu2026}.  Both
branches are part of the input simultaneously.
\end{definition}

\begin{theorem}[Conditional $p$-part identity in rank at most one]\label{thm:closure:t-cascade-rank-le-1}
Let $E/\Q$ have analytic rank $r\in\{0,1\}$ and let $p$ be a prime.
Assume the hypotheses of one of the two branches:
\begin{itemize}
\item[(CM)] $E$ has complex multiplication by $K$; either $p$ is good
  ordinary and split in $K$, or $p\ge5$ is good supersingular and inert in
  $K$; $\Sha(E/\Q)[p^\infty]$ is finite, the Cassels--Tate pairing is
  nondegenerate, the Selmer-complex hypotheses of Sakamoto and Macias
  Castillo--Sano hold, $p\nmid\Tam(E)$, the standard ramification and
  Galois-image hypotheses hold, and a Heegner point is given when $r=1$;
\item[(non-CM)] $E$ has no complex multiplication, the residual
  representation $\bar\rho_{E,p}$ is surjective~\citep{Serre1972}, $p$ is
  good ordinary with $p\nmid N$ and $p\nmid\Tam(E)$; the local,
  residual-irreducibility, ramification and level hypotheses of
  Skinner--Urban hold; $\Sha(E/\Q)[p^\infty]$ is finite, the Cassels--Tate
  pairing is nondegenerate, the Selmer-complex hypotheses hold, and, when
  $r=1$, the Gross--Zagier--Kolyvagin hypotheses hold for a Heegner field.
\end{itemize}
Assume the rank-at-most-one input, and assume further that, under these
branch hypotheses, the results it lists yield the valuation identity
below.  Then
\[
  v_p\!\left(\frac{L^{(r)}(E,1)/r!}{\Omega_E\,\Reg_{\mathrm{NT}}(E)}\right)
  =
  v_p\!\left(\frac{\#\Sha(E/\Q)\,\Tam(E)}{\#E(\Q)_{\tors}^{\,2}}\right).
\]%
\footnote{The formalization records this as a conditional schema.  Lean
proves that the comparison hypothesis, given the forty asserted
propositions and a curve and prime, is equivalent to the bare implication
from the branch hypotheses to the valuation identity; the theorem is
therefore exactly as strong as that implication.}
\end{theorem}

\begin{proof}
Apply the assumed implication to the branch hypotheses.
\end{proof}

The comparison hypothesis in the theorem is where the arithmetic is: in
the non-CM branch, for instance, one divisibility comes from the
Skinner--Urban main conjecture and the other from Euler systems and the
Gross--Zagier--Kolyvagin method, and making the two match up to $p$-adic
units requires the finiteness, nondegeneracy and nondivisibility
hypotheses.  We have not verified that the listed results yield the
identity under exactly the stated hypotheses, and the input is not
organized curve by curve.  The theorem should be read as a statement of
the expected form of the known $p$-part results, not as a new proof of
them, and it does not supply the full identity required by the
recognition predicate in rank at most one.
